\section{复算清单与依赖表}
\subsection{三道短题可分辨哪里出了偏差}
\textbf{题一：三次取整是否同一件事？}
固定一个受服务 $p$。轨道甲两步各欠 $0.2$，轨道乙两步各欠 $0.6$。先在各轨道内部按步求和再上取整，分别需要补 $1$ 与 $2$ 重。若在每步上取整，得到的是 $2$ 与 $2$ 重，会付出更多；若把两个轨道合成一个再取整，只得到 $2$ 重，不能用这个总数同时满足两项独立的轨道约束。最后换算到某个 $w$ 的总下界后才由 $q_w$ 向下取整。注意最后一步甚至可能得到负数，表示收缩原格，不再是在补哪条局部轨道。

\textbf{题二：共用坐标变换哪里真正有用？}
四个源都为一维，目标在数域上有四维。假如四个算术像都恰好落在第一根坐标轴，把每个字的目标端各自转到不同坐标轴，就会让辅助矩阵看似张成四维，而真实像的和只有一维。若整个族共用同一个可逆 $Z$，则
\[
      Z\Bigl(\sum_u \operatorname{im}B_u\Bigr)
      =\sum_u\operatorname{im}(ZB_u),
\]
总维数保持。源侧可以每字不同，因为可逆源换基不改该字的像。实际证明把这一点进一步用于过滤后的相应子空间，然后用真行投影将辅助行扣掉。故共享目标侧换基不是一个格式要求。

\textbf{题三：为何不能把对偶格换成一个同秩格？}
在前面的 $\Zz[i]$ 秩一例子里，$C$ 和 $\pi^{-1}C$ 秩相同，但默认实度量下斜率相差 $\log5$。如果一条不等式给 $s\ge G_a+\mu(C)-M-E_1$，另一条只给
$s\le\mu(\pi^{-1}C)+E_2$，就只能推出
$G_a\le M+E_1+E_2+\log5$，而不能把多出的 $\log5$ 消去。这个例子不说源码真的换错；它告诉审读者该在拼接处比较格的身份和度量，而非只比较秩。

\subsection{四个尚须用传统数学语言展开的接口}
\textbf{局部算术。} 列下界的基准分支和扭转分支应分开书写，说明由哪个等价关系把局部代数列落回到实际配对。正文已将生成元补偿与全局理想贴合分开，但没有从局部 theta 核逐边重算各类素数的端点估值。用同一有理 $p$ 比较辅助域处与系数域处时，还须明列各自赋值单位。\src{Formalism/SharpRestriction/GrowthBounds/WeilJacquetShalikaLaw.lean}{L19-L50}{列界的对象、分支和例外}。

\textbf{周期的非对角界。} 在矩阵层，核的恒等式只告诉读者应该估谁，不能自己给出小量。将异字展开为各类项，解释为什么相同路径可进主项、真正异字需要消失或衰减，以及例外计数如何叠入余项，是更长书稿应做的工作。本章用 LinkS 给定的数值界推到中心平方迹，并未把核的算术生成与分析估计全盘重证。

\textbf{主项高度。} 消项所需的代数恒等式很短，但它从 $S_{\mathrm{Div}}$ 中提取的是已证明的高度主项。完整篇幅不能仅说“恰好抵消”，应写出除法高度、水平点计数及指数成本各自的来源，给同一个族点的参数对照，再算出该恒等式。局部与高度两条界都调用了 $F_o$ 的同一个定义；来源相同不意味着把一条界当作另一条重新使用。

\textbf{最终量词。} 即使中段可以完全独立阅读，仍要从几种窗口和它们的成本比较，到所需解析估计以及目标零点结论逐级交接。最清楚的写法是先写固定精度，再写阈值依赖、在何处取平均、哪些大尺度条件必须满足；最后才使用“任意精度”。\src{Development/AutomorphicArithmetic/PrimitiveCosets/CoreGm.lean}{L12-L22}{后续解析比较如何陈述条件}。

\subsection{源码审读时的身份索引}
从实际端点到最终不等式，有几处即使名称相同也不能只靠拼写判断：
\begin{itemize}
\item $p$ 与 $w$：有理素数及其上方的地点。支持集、付出的幂、以及整数赋值中的贡献要分开。
\item $m,b_u,d_u$：迭代步数、剪前复源维数、实际源格秩。Frame 平均用第三个；最早块界用第二个；局部合计常用第一个。
\item $D,D',D_\star$：全行、真行、算术归一化尺度。三个量通过已证明的比较交接。
\item $\alpha,\mathrm{de},\delta$：实际过滤删掉的源秩比例、建模时剪掉的源秩比例、目标辅助行比例。它们不是一次筛选的三个名字。
\item $B_k,B_u,P_uV_u$：分别是数域格上的映射、实际复配对、全行矩阵。它们的身份通过复嵌入、坐标变换和真行投影联系，不是一个声明本身。
\item $G_a,\mathrm{ASK},\mathrm{Main}$：局部有效收益、待比较的量与主项。$\mathrm{ASK}-G_a$ 的控制在斜率合流后仍有任务。
\end{itemize}
一个可持续的人工审读流程，应让每个新引理能回到这张表上的至少一个具体接口。若某次压缩删了定义但没有留下足够的信息辨认输入，省下的页面要在返工时重新付出。
